\chapter{}

Ghìsola non aveva più il buon contegno di prima. Ambiziosa e caparbia, voleva fare il comodo suo.

Tutte le domeniche, dopo pranzo, fuggiva da casa; e la rivedevano a buio. La nonna andava a cercarla per i poderi: era stata a zonzo per Siena, invece; e, per le strade, le facevano complimenti osceni e proposte di amorazzi. C’era qualcuno che la riconosceva e la seguiva per fermarla e parlarci. Ella sorrideva, un poco stordita e lusingata; perchè non eran contadini ma giovini operai vestiti bene. Quando arrivava alla Porta Camollia, doveva far presto; perchè le guardie daziarie se la mettevano in mezzo e le impedivano di passare.

E quando aveva un fiore, non doveva andare rasente il muro perchè parecchi, ritti su l’uscio delle loro botteghe, allungavano le mani per levarglielo.

Tornata, per non udire brontolii, passava dalla finestra di camera, attaccandosi ai sostegni del pollaio; si spogliava ed entrava a letto senza cenare; arrabbiandosi con il rumore della zuppiera, dove Giacco e Masa mangiavano con i loro cucchiai d’ottone; e quando si sbattevano insieme, Giacco dava un’occhiata a Masa.

Alla fine, la nonna capiva che era in casa; e, pensando che si sarebbe ammalata, le portava di nascosto un pezzo di pane; ma, prima di darglielo, glielo batteva sul capo.

Ghìsola masticava, tenendo il capo volto dalla parte del muro; meravigliandosi che il pane fosse bagnato di lacrime, che non volevano smettere, avendo avuto, poco avanti, piuttosto voglia di ridere. Doveva esser quella la sua vita?

Ma, al rumore dei nonni quando entravano, chiudeva gli occhi; per far credere che dormisse e per il bisogno di non vederli.

L'ultimo giorno che stette a Poggio a' Meli, mentr'era per addormentarsi con una forcella in bocca, che aveva mangiucchiata con i denti, le parve di cadere da una grande altezza e battere sul tetto della casa a Radda: gemendo, si scosse tutta. Il nonno, dall'altro letto, le gridò:

— Stai zitta! Credi che non mi dispiaccia?

Temette d'esser brontolata. Poi riflettè, e a lei parve a voce alta: «Non ci pensano più. Bisogna che non russi».

Ma le dava fastidio l'odore delle lenzuola poco pulite; e, per non sentirlo, se le avvoltolò al collo.

I suoi capelli, sciolti, finivano a punta; e, sopra il capezzale, assomigliavano a una falce.

Le pareva d'entrare in casa: la mamma aveva un vestito nuovo, le due sorelle erano ingrassate. Una voce le chiese:

— Che cosa ci fai qui?

Ed ella rispose:

— Non lo so: non ci sono venuta da me. Ma il babbo dov'è nascosto?

— La colpa è tua.

Ripigliava la voce.

La mamma e le sorelle ascoltavano e guardavano, con un silenzio così orribile ch'ella si slanciava addosso a loro; perchè andassero nell'altra stanza. Ma le pareva di non poter muovere le braccia, e di urtare con il capo in una parete invisibile. Allora sentiva che il cuore cambiava di posto, il ventre faceva lo stesso, la gola si spellava; e i volti della mamma e delle sorelle doventavano spaventevoli. Ella disse:

— Parlate!

Quelle si volsero ad un uscio; e il babbo, con due sacchi pieni su le spalle, con il viso grondante di sangue, tanto sangue che andava a empire la gora del mulino, salì le scale.

Ella, sentendo il peso dei sacchi addosso, urlò.
